\clearpage
\fancyhead[L]{\small CONNAISSANCE DES TEMS.}
\hypertarget{native-497-lower}{}\label{native-497-lower}
\begin{center}
\small Göttingische gelehrte Anzeigen. Nos. 56. 57. pp. 553..562. 7 April 1808.\\[3pt]
\rule{.72\linewidth}{.3pt}
\end{center}

\emph{Connoissance des tems, ou des mouvemens célestes à l'usage des astronomes et des navigateurs pour l'an} 1808, publiée par le bureau des longitudes. Paris. De l'imprimerie impériale. 1806. 502 pages in octavo. --- The same for the year 1809. Ibid. 1807. 502 pages in octavo.

With these volumes of the rightly esteemed astronomical ephemerides, the usual calendar is back in use, after 12 annual volumes (from year IV to XII) appeared in the form of the now forgotten republican calendar. The arrangement of the astronomical calendar itself and of its regular articles, which is retained here unchanged, is too well known for us to dwell on it. We have one observation about the table giving the geographical positions of the principal places on Earth --- which usually receives some additions and improvements from time to time, although even well-known
\SourceLeaf{498}
determinations have yet to be used everywhere, and some gross misprints are carried over from one volume to the next --- with the year 1809 it is now arranged in a single alphabetical sequence for the whole Earth, whereas it was previously arranged by country. The reviewer would prefer the previous arrangement to this innovation, since it is sometimes convenient to have the reliable determinations of positions for a given country together. We therefore turn at once to the \emph{supplements}, which always comprise a great number of valuable papers, observations and interesting notes, and give the ephemerides lasting value for astronomers.

The editing of the supplements, previously done by \textsc{Lalande}, was taken over by \textsc{Delambre} with the year 1808. In that year the editors also began to open them with the observations made by \textsc{Bouvard} at the Imperial Observatory in Paris and will continue to do so. This regular publication of observations as numerous as they are valuable, in the same manner as the Greenwich observations, must be very welcome to astronomers and considerably increases the value of the Connoissance des tems. In the 1808 volume, after a short description of the instruments, the observations from 24 September 1803 to 21 September 1804 take up 83 pages. The excellent transit instrument by \textsc{Berge} has a focal length of $2\frac{1}{2}$ metres (92 Paris inches) and an aperture of 11 centimetres (4 inches). It was installed in September 1803. The clock is by \textsc{Louis Berthoud}. The eight-foot mural quadrant by \textsc{Bird}, fitted with an achromatic telescope, is the one that formerly belonged to \textsc{Lemonnier}. Among the observations we also find numerous observations of Ceres and Pallas that have yet to be used. Observations and parabolic elements of the comet of 1804 and of the two of 1805 are appended, though for the latter only the reduced longitudes and latitudes are given. This astronomical journal by \textsc{Bouvard} is followed by observations by \textsc{Flaugerques} at Viviers and \textsc{Vidal} at Toulouse. The latter consist of culminations of Mercury and Venus, some observed by this sharp-eyed astronomer in the same minute as the Sun. Solar eclipses and occultations of planets and stars, observed by \textsc{Scarpellizi} at Rome and calculated in part by \textsc{Lalande}. From \textsc{Lalande}, observations of the inferior conjunctions of Venus in 1802 and 1804 and of the solar eclipse of 1 February 1804 from several places. From \emph{the very same astronomer}, a new catalogue of the proper motions of 500 stars. The reviewer counts 28 in which the annual proper motion in right ascension exceeds half a second of arc, and 15 in which it does so in declination. But it seems that \textsc{Lalande} has paid insufficient attention to the nonuniformity of precession, although he expressly says that he has taken it into account. At least we find that he did not do so for the pole star. From \textsc{Lacaille}’s position for 1750 and \textsc{von Zach}’s for 1800, the observed annual motion in right ascension is simply inferred to be $173''18$, and is compared with the motion calculated for 1800 (instead of 1775, as it should be), $207''52$. The annual proper motion is then quite wrongly put at $34''34$. --- From \textsc{Prony}, a paper on calculating geographical longitude and latitude from distances from the meridian and its perpendicular on the elliptical spheroid. His formulas apply only to moderate distances, but their simplicity recommends them. From \textsc{Delambre}, a new and ingenious method for finding the configuration of Jupiter’s satellites by calculation alone, using certain auxiliary tables. Mechanical aids had generally been used for this purpose until now. From \emph{the same astronomer}, the history of astronomy for 1804 and 1805. The discovery of \emph{Juno} rightly takes first place here. The judgement that any astronomer would always count himself fortunate to buy so great a discovery at the cost of a whole year --- before the discovery of Ceres, surely no one would have hesitated to say a whole lifetime --- of strenuous labour does credit to the French
\SourceLeaf{499}
astronomer. The fact that astronomy has had the good fortune to gain so priceless an addition four times in six years can hardly diminish its value. On the contrary, these events, so glorious chiefly for Germany, must spur on astronomers’ courage and activity all the more, since they so greatly increase the hope of similar successes in the future. All the more surprising, and wholly at odds with that judgement, is the coldness with which \textsc{Delambre} speaks of the plan drawn up seven years ago in Lilienthal, before Ceres was discovered, chiefly by Mr \textsc{von Zach}, to divide the entire sky among a number of astronomers for continued searches. He too was invited to take part. He considers such a comparison extremely laborious and boring, requiring one to stand mechanically at a stationary telescope every night for hours, counting stars. He believes it would put off altogether anyone who felt a calling for astronomy. One should merely wait for a lucky chance \emph{in passing}, so to speak, whenever one observes the Moon, the other planets and the larger fixed stars. The only person who \emph{deserves} to make such a discovery is someone who happens to make it without setting out to do so, while engaged in some great and \emph{truly useful} undertaking, and so on. Statements of this kind from a man who deserves so much respect for his many distinguished works cannot be passed over in silence. Men. whose names carry authority ought, it seems to us, to be all the more careful about expressing judgements that might seem to reflect the spirit of our narrow-minded and frivolous century, in which people think they honour the Muses by reducing them to mere servants of our needs, or in which, as one of our wittiest writers puts it, the stars in the most sublime astronomy are valued only as step counters and signposts for fleets carrying pepper. To anyone who has too little enthusiasm for science to feel the greatness of discovering a previously unknown celestial body that, like the Earth, regularly circles our Sun and whose appearances will still occupy observers of the heavens thousands of years from now, and whose first thought is what use such a discovery will be, we have only our conviction to offer. These new planets themselves will give a wholly new impetus to theoretical and physical astronomy. But anyone who takes a different view and sees in the splendid success of our \textsc{Olbers} and \textsc{Harding} a reason for the well-founded hope that more successes of the same kind can be expected, must earnestly wish to see that plan put into effect soon. It alone can make us independent of chance in an undertaking whose scope is manageable for individual astronomers. Such a reader will be very glad to hear that an important step has already been taken. It is with real pleasure that we report that the sky charts of our Professor \textsc{Harding}, which have laid such excellent groundwork for the undertaking, are now very close to publication despite the difficulties the sciences face here in our times. In any case, \textsc{Delambre} has no grounds for thinking that searches of this kind are so arduous, time-consuming and tedious. With a suitable method, using good night telescopes with strong light-gathering power instead of fixed instruments, it is possible to survey a considerable area in a short time once complete charts are available. One will surely always find this occupation a pleasant respite from other mechanical tasks of observational and computational astronomy. Nor is it at all necessary to repeat the whole survey every night.

After this digression, which seemed to us a timely word, we return to the remaining articles of that paper. On the comets of 1804 and 1805. \textsc{Schröter}’s and \textsc{Harding}’s
\SourceLeaf{500} observations of Saturn’s ring. \textsc{Schröter}’s and \textsc{Herschel}’s determinations of the diameters of the new planets. \textsc{Piazzi}’s observations of the parallax of some fixed stars. Some formulas by \textsc{Camerer} and \textsc{Olbers} borrowed from the Berlin astronomical yearbooks. On \textsc{Bowditch}’s and \textsc{Mendoza}’s methods for reducing apparent lunar distances to true ones. Reports from the second volume of the \emph{Asiatic Researches} on the Hindus’ methods for calculating sine tables. The remaining shorter articles in this volume contain notices of various new astronomical works, the official proceedings on the abolition of the republican calendar and the reintroduction of the Gregorian calendar, observations of the solar eclipse of 16 June 1806 and the latest transit of Mercury, the announcement concerning the prize offered for calculating the perturbations of Pallas, some improvements to the tables in \textsc{Lalande}’s astronomy, and an extract from the meteorological observations at Paris for year XII.

In the \emph{volume} for 1809, the observations made at the Imperial Observatory from 23 September 1804 to the end of 1805 take up 108 pages. The new planets, Ceres, Pallas and Juno, were diligently observed at the meridian in 1804, the last of them until 24 December. We are sorry, however, to find no observations of Pallas and Ceres from the end of 1805. For the two comets of 1805 we find the original observations here. There follow five years of astronomical observations by \textsc{Messier}, from 1760\mbox{. . . .}1764, continuing the observations already published in earlier volumes of the Connoissance des tems. Here one finds a solar eclipse, two lunar eclipses, several occultations of fixed stars, one of Mars as well, and a great many eclipses of Jupiter’s satellites. For enthusiasts of meteorology we single out the remark that on 19 July 1760 between four and five in the afternoon the mercury thermometer in Paris rose to 32 degrees. --- Astronomical observations at Viviers by \textsc{Flaugerques}. --- From \textsc{Burckhardt}, the arrangement of the gearing for representing the motions of the planets, some remarks on making the use of \textsc{Borda}’s circle more precise, general and convenient on various occasions, and investigations of the proper motion of our solar system by comparing the observed proper motions of the fixed stars. The result is that little can yet be determined about it with any satisfaction. Some older Chinese observations of solstices from a manuscript formerly sent by Father \textsc{Gaubil} to \textsc{Delisle} and kept in the Imperial Library. Also from \textsc{Burckhardt}, a proposal for a new design for reflecting telescopes. At half the focal length he intercepts the rays reflected by the large mirror with a flat mirror perpendicular to the axis, which sends them back to the eyepieces in the pierced concave mirror. The advantage over \textsc{Newton}’s design is that with the same focal length for the mirror, the whole telescope is only half as long. On the other hand, the small mirror (which must be half as large in diameter) takes from the large mirror a quarter of the light entering the tube. The instrument maker \textsc{Caroché} will try to carry out the proposal. --- Remarks on a method proposed by \textsc{Ducum} at Bordeaux for determining latitude at sea from two altitudes measured off the meridian. The writer of the account has added several observations of his own on this problem and on \textsc{Douwes}’s solution. --- From \textsc{Oltmanns}, a valuable investigation of the longitude of Quito using lunar distances and eclipses of Jupiter’s satellites observed by Mr \textsc{v. Humboldt}, and various older observations. The result he finds for this longitude, for which determinations have varied so widely until now, is $5^{\mathrm u}\,24^{\mathrm m}\,20^{\mathrm s}$ west of Paris. --- On a passage in Ptolemy’s Geography concerning the projection used by Marinus of Tyre, in which people had quite wrongly thought they could already detect a trace of Mercator’s projection. --- Observations of Ceres, Pallas and Juno by \textsc{Poczobut} and \textsc{Keschka} at Wilna. --- From \textsc{Henry}, exact formulas for the parallax in longitude and
\SourceLeaf{501} latitude. --- A new kind of micrometer by \textsc{Prony}. The underlying idea is as simple as it is ingenious and can also be put to use on other occasions. To give the movable wire (the slider) a very slow motion without needing a screw with very fine threads, which is difficult to make and soon changes, he drives it with a screw with fairly coarse threads. On both sides of these threads, however, he cuts further threads on the same axis, with the same pitch on both sides and a pitch slightly smaller or larger than that of the middle threads. Their nuts are fixed in the housing. In this way the actual motion of the slider depends only on the difference between the two kinds of thread. --- Extracts from some new astronomical writings. The extracts from a paper by \textsc{Monteiro} on the calculation of eclipses in the fourth volume of the Coimbra ephemerides are of interest. --- The history of astronomy for 1807 concerns mainly the discovery of Vesta. German readers will find nothing new about it here. For the latest comet we also find here the first observations by \textsc{Thulis} on 21 and 22 September, with the declinations given only in whole minutes. --- Some geographical determinations in the Mediterranean Sea and on its coast, by \textsc{Dionisio Alcala Galeano}. --- This volume again concludes with an extract from \textsc{Bouvard}’s meteorological observations for year XIII.
